\documentclass[11pt,a4paper]{article}
% =====================================================================
%  Shared preamble for the Part V lecture notes (lectures 19-22).
%
%  These four lectures sit outside Geron, so the notes ARE the primary
%  source and are examined on the same terms as the textbook parts.
%  Everything here is chosen to make that readable on paper: one column,
%  generous measure, and the same six-colour palette as the slides so a
%  student moving between the two is not re-learning what red means.
%
%  Build with pdflatex (twice, for the toc and the refs):
%      cd notes && latexmk -pdf lecture-19.tex
%  or  make -C notes
% =====================================================================
\usepackage[T1]{fontenc}
\usepackage[utf8]{inputenc}
\usepackage{newtxtext,newtxmath}      % Times-like: prints well, reads well
\usepackage[scaled=0.86]{inconsolata} % code, at a size that sits beside Times
\usepackage{microtype}
% newtx defines \Bbbk, \openbox, \proof and \endproof; amssymb and amsthm
% then redefine them and abort. Freeing the four names before those two
% packages load is the standard fix, and it costs nothing here: we use
% amsthm only for \newtheorem.
\let\Bbbk\relax
\usepackage{amsmath,amssymb}
\let\openbox\relax
\let\proof\relax
\let\endproof\relax
\usepackage{amsthm}
\usepackage{booktabs}
\usepackage{enumitem}
\usepackage{xcolor}
\usepackage{tcolorbox}
\usepackage{titlesec}
\usepackage{graphicx}
\usepackage[hidelinks]{hyperref}
\tcbuselibrary{skins,breakable}

% ---- the course palette, from assets/css/custom.css -------------------
\definecolor{aimlprimary}{HTML}{0B3D62}   % deep blue  - structure
\definecolor{aimlaccent} {HTML}{C0392B}   % brick red  - failures
\definecolor{aimlsuccess}{HTML}{14663A}   % green      - repairs
\definecolor{aimlmath}   {HTML}{6C3483}   % purple     - the derivation
\definecolor{aimlmuted}  {HTML}{4B5563}
\definecolor{aimlrule}   {HTML}{B0BCC7}
\definecolor{aimlsoft}   {HTML}{F4F7F9}

\usepackage[a4paper,margin=32mm,top=30mm,bottom=30mm]{geometry}
\setlength{\parskip}{0.45\baselineskip}
\setlength{\parindent}{0pt}
\linespread{1.06}

% ---- headings --------------------------------------------------------
\titleformat{\section}{\Large\bfseries\color{aimlprimary}}{\thesection}{0.7em}{}
\titleformat{\subsection}{\large\bfseries\color{aimlprimary}}{\thesubsection}{0.7em}{}
\titleformat{\subsubsection}{\bfseries\color{aimlmuted}}{}{0em}{}
\titlespacing*{\section}{0pt}{1.6\baselineskip}{0.5\baselineskip}
\titlespacing*{\subsection}{0pt}{1.2\baselineskip}{0.35\baselineskip}

% ---- boxes -----------------------------------------------------------
% One box per role, and the role is the colour. A student who has seen the
% slides already knows what each one means before reading a word of it.
\newtcolorbox{keybox}[1][]{
  breakable, enhanced, colback=aimlsoft, colframe=aimlprimary,
  boxrule=0.4pt, left=8pt, right=8pt, top=6pt, bottom=6pt,
  fonttitle=\bfseries\small, coltitle=white, colbacktitle=aimlprimary,
  attach boxed title to top left={yshift=-2mm, xshift=4mm},
  boxed title style={boxrule=0pt, arc=1pt}, #1}

\newtcolorbox{failbox}[1][]{
  breakable, enhanced, colback=white, colframe=aimlaccent,
  boxrule=0.9pt, left=8pt, right=8pt, top=6pt, bottom=6pt,
  fonttitle=\bfseries\small, coltitle=white, colbacktitle=aimlaccent,
  attach boxed title to top left={yshift=-2mm, xshift=4mm},
  boxed title style={boxrule=0pt, arc=1pt}, #1}

\newtcolorbox{derivbox}[1][]{
  breakable, enhanced, colback=white, colframe=aimlmath,
  boxrule=0.6pt, left=8pt, right=8pt, top=6pt, bottom=6pt,
  fonttitle=\bfseries\small, coltitle=white, colbacktitle=aimlmath,
  attach boxed title to top left={yshift=-2mm, xshift=4mm},
  boxed title style={boxrule=0pt, arc=1pt}, #1}

% An exercise. Deliberately the same shape as the notebook's `try` field:
% one change, and what should happen to the output.
\newtcolorbox{trybox}[1][]{
  breakable, enhanced, colback=white, colframe=aimlsuccess,
  boxrule=0.5pt, left=8pt, right=8pt, top=6pt, bottom=6pt,
  fonttitle=\bfseries\small, coltitle=white, colbacktitle=aimlsuccess,
  attach boxed title to top left={yshift=-2mm, xshift=4mm},
  boxed title style={boxrule=0pt, arc=1pt}, #1}

% ---- small semantics -------------------------------------------------
\newcommand{\scope}[1]{\textcolor{aimlmuted}{\small #1}}
\newcommand{\term}[1]{\textbf{#1}}
\newcommand{\code}[1]{\texttt{\small #1}}
\newcommand{\lecture}[1]{Lecture~#1}

\newtheorem{definition}{Definition}[section]
\newtheorem{proposition}[definition]{Proposition}

% ---- title block -----------------------------------------------------
% \notestitle{19}{Information retrieval: the lexical foundation}{SciFact (BEIR)}
\newcommand{\notestitle}[3]{%
  \thispagestyle{empty}
  \begin{flushleft}
    {\color{aimlmuted}\small\sffamily
     APPLICATIONS OF MACHINE LEARNING \quad$\cdot$\quad
     BSc Mathematics of Artificial Intelligence}\\[2mm]
    {\color{aimlrule}\rule{\textwidth}{0.8pt}}\\[4mm]
    {\color{aimlmuted}\large Lecture #1 \quad$\cdot$\quad Part V \quad$\cdot$\quad
     Extended notes}\\[3mm]
    {\Huge\bfseries\color{aimlprimary}\baselineskip=1.15\baselineskip #2\par}\vspace{4mm}
    {\color{aimlmuted}\large Dataset: #3}\\[3mm]
    {\color{aimlrule}\rule{\textwidth}{0.8pt}}
  \end{flushleft}
  \vspace{2mm}
  \begin{keybox}[title={Why these notes exist}]
    Lectures 19--22 sit outside G\'eron, the course's primary text. For those
    four lectures \emph{these notes are the primary source}, and they are
    examined on exactly the same terms as the textbook parts. Everything in
    the slide deck for this lecture is here, with the arguments written out at
    length rather than compressed onto a slide; the numbers are the ones the
    accompanying notebook prints, and every one of them is a measurement on
    the corpus named above rather than a figure quoted from a paper.
  \end{keybox}
  \vspace{1mm}
  {\small\tableofcontents}
  \vspace{2mm}
  \noindent{\color{aimlrule}\rule{\textwidth}{0.4pt}}
  \vspace{1mm}
}


\begin{document}
\notestitle{20}{Information retrieval: dense retrieval}{SciFact (BEIR), the same 5{,}183 abstracts}

\section{What an index of strings cannot do}

\lecture{19} ended with a measurement rather than an opinion: a relevant
abstract is missing 40.9\% of its claim's words on average, and 28.6\% of
judged pairs share fewer than half. \emph{Heart attack} and \emph{myocardial
infarction}; \emph{lower numbers} and \emph{reduced inoculum}; a mechanism
described rather than named.

\begin{keybox}[title={The requirement}]
A representation in which two texts that mean the same thing are \emph{close},
whether or not they share a word. Which is precisely what Lectures 17 and 18
built --- and the reason Part~V comes after Part~IV rather than before it.
\end{keybox}

\subsection{The three repairs that do not work}

\begin{center}
\small
\begin{tabular}{p{0.28\textwidth}p{0.62\textwidth}}
\toprule
\textbf{Repair} & \textbf{Why it falls short} \\
\midrule
Stemming & merges word forms, not concepts --- nothing links
\emph{infarction} to \emph{attack} \\
A synonym dictionary & someone must write it, per domain, and it cannot cover
paraphrase \\
Query expansion & adds terms from the top results, so it needs the answer to
already be near the top --- which is exactly what has failed \\
\bottomrule
\end{tabular}
\end{center}

Each is a patch on the same assumption: that relevance is visible in shared
strings. The assumption is what has to go. This is the shape of every move from
feature engineering to learned representations in this course --- the patches
are not stupid, they are \emph{bounded}, and the bound is the assumption.

\subsection{Be fair to the lexical index first}

Before replacing BM25, note what it is exactly right about.

\begin{itemize}[leftmargin=*,nosep]
  \item A rare technical term --- a gene name, a drug, an identifier --- is
        matched \emph{exactly}, where an embedding may blur it into its
        neighbours.
  \item It needs no training data, no accelerator, and no model version to
        record.
  \item Its score is explainable: the contribution of every term can be
        printed, as we did in \lecture{19}.
  \item A new document is indexed in microseconds. No re-encoding, no
        re-fitting.
\end{itemize}

Three of those four are operational requirements that outlive whichever model
is fashionable, and they are why BM25 is still in the first stage of systems
that also run neural rankers.

\section{Bi-encoders}

Part~IV already gave us everything but one idea. \lecture{17} gave a token
embedding and a sentence vector from pooling; \lecture{18} gave attention, so
that a token's vector depends on its context. Both give a fixed-width vector
for a variable-length text.

\begin{keybox}[title={The one new idea}]
Put the query and the document in the \emph{same space}, and make the score a
dot product. Everything else in this lecture is a consequence of that choice:
what can be precomputed, what must be approximated, and what the model is
allowed to see.
\end{keybox}

There is no separate derivation in this lecture for that reason. The
mathematics is a normalised inner product, and you have had it since
\lecture{8}.

\subsection{From tokens to one vector}

A retrieval score needs one vector per \emph{text}, so the token vectors must be
pooled.

\begin{center}
\small
\begin{tabular}{p{0.16\textwidth}p{0.72\textwidth}}
\toprule
\textbf{Pooling} & \textbf{What it does} \\
\midrule
\code{[CLS]} & take the reserved first token's vector --- meaningful only if
the model was trained to put something there \\
Mean & average the token vectors, \emph{masking padding}; robust, and what our
model uses \\
Max & component-wise maximum; rarely better, and discards how often a thing was
said \\
\bottomrule
\end{tabular}
\end{center}

\begin{failbox}[title={Masking the padding is not optional}]
Mean-pooling over the padding tokens makes a text's vector depend on
\emph{the batch it was encoded in}. The same abstract gets two different
vectors, the retrieval scores change, and nothing raises an error.
\end{failbox}

\subsection{Why the vectors are normalised}

Normalise to unit length and the inner product becomes the cosine, which
discards magnitude. That is deliberate.

\begin{itemize}[leftmargin=*,nosep]
  \item Magnitude in these encoders tracks length and token frequency, not
        relevance --- so an unnormalised dot product quietly prefers long
        documents, which is the failure BM25's $b$ was invented for.
  \item Normalised scores are bounded in $[-1,1]$, which makes a threshold
        transferable between corpora.
  \item And it makes the clustering of Section~\ref{sec:ann} legitimate:
        $k$-means on the sphere is a meaningful thing to do.
\end{itemize}

The same trap as \lecture{19}'s length normalisation, arriving from a different
direction. Long documents win by default unless something is done about it.

\subsection{Encode both sides into one space}

\[
  \text{score}(q,d) \;=\; E(q)\cdot E(d),
  \qquad E : \text{text} \to \mathbb{R}^{384}.
\]

One encoder, applied \emph{separately} to the query and to the document.

\begin{keybox}[title={Why ``separately'' is the whole design}]
$E(d)$ does not depend on the query. So every document can be encoded once,
offline, and \emph{stored}; at request time you encode one short query and take
5{,}183 inner products. The expensive part happens before anyone asks.

Hold on to that sentence. It is also the reason the bi-encoder is the weaker of
the two neural models in this lecture.
\end{keybox}

\subsubsection{One encoder or two?}

\begin{center}
\small
\begin{tabular}{p{0.17\textwidth}p{0.33\textwidth}p{0.33\textwidth}}
\toprule
 & \textbf{Shared encoder} & \textbf{Two encoders} \\
\midrule
Parameters & one set & two sets \\
Suits & symmetric tasks --- sentence similarity & asymmetric --- short query,
long document \\
Risk & a short query and a long abstract are not the same kind of text & twice
the parameters on the same labels \\
\bottomrule
\end{tabular}
\end{center}

We use a shared encoder because the model we borrow was trained that way. It is
a real modelling choice and it should be reported: a paper that says
``bi-encoder'' without saying which has told you less than it thinks.

\subsection{What it costs to store}

\begin{center}
\small
\begin{tabular}{lrr}
\toprule
 & \textbf{Per document} & \textbf{Whole corpus} \\
\midrule
Dense vectors, 384-d float32 & 1{,}536 bytes & 8.0 MB \\
Inverted index postings      & 488 bytes    & 2.5 MB \\
\bottomrule
\end{tabular}
\end{center}

Three times the storage --- and, unlike postings, \emph{every byte of it is
touched by every query}. The inverted index skips the documents sharing no
term; the dense index cannot skip anything, because every vector has a non-zero
dot product with every query. That is why approximate nearest-neighbour search
exists, and it is Section~\ref{sec:ann}.

The dimension is a dial too:

\begin{center}
\small
\begin{tabular}{lr}
\toprule
Bytes per document                        & 1{,}536 \\
This corpus                               & 8.0 MB \\
The same corpus at 768 dimensions         & 15.9 MB \\
Ten million documents at 384-d            & 15.4 GB \\
\bottomrule
\end{tabular}
\end{center}

A dense index is expensive not to build but \emph{to use}, and that is what
quantisation and approximate search exist to reduce.

\subsection{The whole retrieval system, in six lines}

\begin{verbatim}
D = encode(model, documents)              # once, offline, 5,183 x 384
D = D / np.linalg.norm(D, axis=1, keepdims=True)

def search(query, k=10):
    q = encode(model, [query])[0]
    q /= np.linalg.norm(q)
    s = D @ q                             # every document, exactly
    return np.argsort(-s)[:k]
\end{verbatim}

Note what is \emph{not} here: no index structure, no library, no database. The
matrix product is the search, and on this corpus it is fast enough to be exact.
Everything a vector database adds is a way of not doing that matrix product in
full, and it is bought with recall.

\section{Training one, and why we will not}

\subsection{What the objective must say}

We want $E(q)\cdot E(d^{+})$ large and $E(q)\cdot E(d^{-})$ small. Written as a
softmax over one positive and a set of negatives:

\begin{derivbox}[title={The contrastive objective}]
\[
  \mathcal{L} \;=\; -\log
  \frac{\exp\!\big(E(q)\cdot E(d^{+})/\tau\big)}
       {\exp\!\big(E(q)\cdot E(d^{+})/\tau\big)
        + \sum_{d^{-}}\exp\!\big(E(q)\cdot E(d^{-})/\tau\big)}
\]
which is the cross-entropy of \lecture{17} with the \emph{documents} as
classes, and the objective you meet again in \lecture{23} with images on one
side.
\end{derivbox}

The positives come from the judgements. The negatives are the problem: there
are millions of documents and you cannot score them all.

\subsection{In-batch negatives}

In a batch of $B$ query--document pairs, every \emph{other} document in the
batch is a negative for every query. One encoding pass gives $B^2$ scores.

\begin{center}
\small
\begin{tabular}{rrrr}
\toprule
\textbf{Batch $B$} & \textbf{Texts encoded} & \textbf{Scores available} &
\textbf{Negatives per query} \\
\midrule
8   & 16    & 64      & 7 \\
32  & 64    & 1{,}024 & 31 \\
128 & 256   & 16{,}384 & 127 \\
512 & 1{,}024 & 262{,}144 & 511 \\
\bottomrule
\end{tabular}
\end{center}

Encoding grows linearly and negatives grow linearly; the \emph{scores} grow
quadratically, because the dot products are nearly free next to the encoder.
This is why dense-retrieval papers report batch sizes in the thousands, and why
the batch size here is a \emph{modelling decision} rather than a memory
setting.

\begin{failbox}[title={``In-batch negatives are free''}]
The scores are free. The encodings are not, and the score matrix itself is not:
at $B = 512$ it is about a megabyte in float32, and at $B = 16{,}384$ it is a
gigabyte. What actually caps the batch is memory, and the batch is the number
of classes in the problem you are solving.
\end{failbox}

\subsection{The temperature}

The $\tau$ in the objective is not decoration. Scores are cosines, so they live
in $[-1,1]$; a softmax over that range is nearly uniform and the gradient
nearly vanishes. Small $\tau$ sharpens the distribution, so the hardest
negative dominates the gradient; large $\tau$ flattens it, so all negatives
contribute almost equally and the model learns slowly.

Typical values are 0.01 to 0.05, which are \emph{small}: dividing a cosine by
0.05 spreads $[-1,1]$ across $[-20,20]$, which is where a softmax has gradient
to give. \lecture{23} derives this properly.

\subsection{Why the negatives are most of the work}

A random negative is almost always trivially wrong: a random abstract shares
nothing with the claim, so the model learns to separate ``about medicine'' from
``about anything else''.

\begin{itemize}[leftmargin=*,nosep]
  \item \term{Hard negatives} --- documents a first stage ranks highly but
        which are not relevant. These teach the distinction that matters.
  \item \term{False negatives} --- and here is the trap: a hard negative mined
        this way is often \emph{relevant but unjudged}, so training actively
        teaches the model that a correct answer is wrong.
\end{itemize}

\begin{failbox}[title={The pooling problem, now inside training}]
\lecture{19} said unjudged documents are scored as irrelevant \emph{at
evaluation}. Here they are scored as irrelevant \emph{during training}, which
is worse: the error is no longer an accounting mistake, it is a gradient.

On this corpus, of 1.55 million claim--abstract pairs only 338 are judged
relevant, a mean of 1.13 per claim. BM25's top ten for a claim contains on
average 0.86 judged-relevant abstracts and 9.14 others --- and those others are
\emph{unjudged}, not judged irrelevant. Train on them as negatives and the
gradient pushes the model away from correct answers. This is the single
commonest reason a fine-tuned retriever underperforms the paper it copied.
\end{failbox}

\subsection{Today we do not train one}

SciFact has 300 test claims and a training set of a few hundred more.
Fine-tuning a bi-encoder on that would measure the split, not the method. So we
take a model trained on something else entirely --- \code{all-MiniLM-L6-v2},
trained by contrastive learning on about a billion general sentence pairs
scraped from the web, never shown a scientific claim --- and apply it
\term{zero shot}.

That is also the honest question. ``Should I put an off-the-shelf embedding
model in front of my search?'' is the decision most of you will actually face.
It is not ``can a dense retriever be trained to beat BM25 given enough
labels?'' --- the answer to that is yes, and it is not usually the question you
have.

\begin{trybox}[title={Predict before you read on}]
A model whose training distribution excludes your domain should do worst
exactly where the domain vocabulary is most specialised. BM25 scored 0.6611
NDCG@10 with no training at all. Write down a number for the zero-shot dense
retriever \emph{now}. The point is not whether you are right; it is to notice,
on the next page, whether you would have argued for the answer had you seen it
first.
\end{trybox}

\subsection{What the bi-encoder gave up}

A single 384-dimensional vector must stand for a 225-token abstract, and it is
computed \emph{before the query is known}.

\begin{itemize}[leftmargin=*,nosep]
  \item An abstract about three things compresses to one point, so it is
        somewhat about each and exactly about none.
  \item Nothing in the encoder can emphasise the part of the document the query
        is asking about --- it has not been asked yet.
  \item Exact term matching is not guaranteed: two different gene names may
        land close together.
\end{itemize}

That is the price of indexability, and the trade is \emph{structural}: it is
not fixed by a bigger encoder. It is the whole reason a second, non-indexable
model appears later in this lecture.

\subsection{The operational cost the accuracy table never shows}

\begin{center}
\small
\begin{tabular}{p{0.2\textwidth}p{0.3\textwidth}p{0.36\textwidth}}
\toprule
 & \textbf{Inverted index} & \textbf{Dense index} \\
\midrule
New document & append to its terms' postings; microseconds & one encoder pass,
then append the vector \\
New model version & nothing to do & re-encode the entire corpus \\
Mixed versions & impossible & silently wrong --- vectors from two models are
not comparable \\
\bottomrule
\end{tabular}
\end{center}

The last row is a real outage, and a quiet one: nothing errors, the scores are
simply meaningless for the documents encoded by the old model. A dense index
needs a model version stored beside it.

\section{Against BM25, on the same judgements}

\subsection{The protocol, stated before the numbers}

Same corpus of 5{,}183 abstracts, same 300 test claims, same judgements. Same
metric code --- the functions written in \lecture{19}'s notebook, unchanged.
Both retrievers scored \emph{exactly}: every document for every query, no
approximate index in either. And the metric was named before the run: NDCG@10,
with recall@100 reported alongside because that is what a first stage is judged
on.

\subsection{The result}

\begin{center}
\small
\begin{tabular}{lrr}
\toprule
\textbf{Metric} & \textbf{BM25} & \textbf{Dense, zero shot} \\
\midrule
NDCG@10     & 0.6611 & 0.6451 \\
Precision@1 & 0.5433 & 0.5033 \\
MRR         & 0.6350 & 0.6113 \\
Recall@10   & 0.7784 & 0.7833 \\
Recall@100  & 0.8852 & 0.9250 \\
\bottomrule
\end{tabular}
\end{center}

\begin{keybox}[title={The neural model loses on the headline metric}]
Thirty years of age, no training, no accelerator, no vector store --- and BM25
is ahead by 0.0160 NDCG@10.
\end{keybox}

It is not a bug in the experiment: same corpus, same claims, same judgements,
same metric code. It is a well-known result --- BEIR was built to report it,
and out-of-domain dense retrieval loses to BM25 on a good fraction of its
tasks. Scientific text is exactly where it should be worst, because the
vocabulary is technical and a general-purpose encoder has never seen most of
it.

\subsubsection{What this result does not license}

\begin{itemize}[leftmargin=*,nosep]
  \item It does \emph{not} say dense retrieval is worse in general. Fine-tuned
        on in-domain pairs it beats BM25 on this corpus, and that is
        well replicated.
  \item It does \emph{not} say embeddings are a bad idea. Recall@100 went up,
        and the mismatch case went from rank 4{,}999 to rank 8.
  \item It \emph{does} say that this particular substitution --- off-the-shelf
        encoder for BM25, no labels, technical corpus --- is a loss, and that
        is a decision many teams make without measuring.
\end{itemize}

Stating the scope of a claim is not hedging. A result whose boundary you cannot
draw is a result you cannot use.

\subsection{But look at recall@100}

\begin{center}
\small
\begin{tabular}{lrrr}
\toprule
\textbf{Metric} & \textbf{BM25} & \textbf{Dense} & \textbf{Difference} \\
\midrule
NDCG@10    & 0.6611 & 0.6451 & $-0.0160$ \\
Recall@10  & 0.7784 & 0.7833 & $+0.0049$ \\
Recall@100 & 0.8852 & 0.9250 & $+0.0398$ \\
\bottomrule
\end{tabular}
\end{center}

The dense retriever is \emph{worse at ordering and better at finding}. Recall
from \lecture{19} what a first stage is judged on: not what the user sees, but
what it makes possible. On that criterion --- the criterion for the job it
would actually be given --- the dense retriever wins.

\subsection{Do they find the same documents?}

Take every relevant abstract in the test set and ask which method placed it in
the top ten.

\begin{center}
\small
\begin{tabular}{lrr}
\toprule
\textbf{Found by} & \textbf{Count} & \textbf{Share} \\
\midrule
Both        & 227 & 67.0\% \\
BM25 only   & 30  & 8.8\% \\
Dense only  & 38  & 11.2\% \\
Neither     & 44  & 13.0\% \\
\bottomrule
\end{tabular}
\end{center}

68 of 339 relevant abstracts are found by exactly one of the two. They are not
two attempts at the same ranking; they fail on \emph{different documents}.

\begin{keybox}[title={A mean hides the disagreement}]
BM25 0.6611 against dense 0.6451 sounds like two systems doing almost the same
thing. Per claim they frequently disagree completely. Two systems with the same
mean can have entirely different error sets: the mean answers ``which is better
on average'' and cannot answer ``are these the same system'' --- and only the
second question tells you whether to combine them.
\end{keybox}

\subsection{The claim BM25 could not find}

\emph{``In mice, P. chabaudi parasites are able to proliferate faster early in
infection when inoculated at lower numbers than when inoculated at high
numbers.''}

\begin{center}
\small
\begin{tabular}{lr}
\toprule
\multicolumn{2}{l}{\textbf{Rank of its one relevant abstract}} \\
\midrule
BM25             & 4{,}999 \\
Dense, zero shot & 8 \\
\bottomrule
\end{tabular}
\end{center}

Nothing about the text changed; the abstract still shares almost no words with
the claim. What changed is that the comparison is no longer between strings. Be
precise about what this shows: not that the dense retriever is better --- the
table above says it is not --- but that the two methods fail
\emph{independently}.

\section{Combining them, and paying for it}

\subsection{Hybrid retrieval}

The scores are not on a comparable scale --- BM25 returns unbounded sums, the
dense model returns cosines --- so fuse the \emph{ranks} instead:

\begin{derivbox}[title={Reciprocal rank fusion}]
\[
  \text{RRF}(d) \;=\; \sum_{\text{systems}} \frac{1}{k + \text{rank}(d)},
  \qquad k = 60 .
\]
No training, no tuning, no calibration, and it needs nothing from the two
systems but their orderings --- which is why it is what people actually deploy.
The constant $k$ damps the top ranks so one system cannot dominate on a single
confident hit; 60 is conventional, in the same sense that $\log_2(i+1)$ is.
\end{derivbox}

\begin{center}
\small
\begin{tabular}{lrrr}
\toprule
\textbf{Metric} & \textbf{BM25} & \textbf{Dense} & \textbf{Hybrid} \\
\midrule
NDCG@10     & 0.6611 & 0.6451 & 0.6948 \\
Precision@1 & 0.5433 & 0.5033 & 0.5633 \\
Recall@10   & 0.7784 & 0.7833 & 0.8239 \\
Recall@100  & 0.8852 & 0.9250 & 0.9617 \\
\bottomrule
\end{tabular}
\end{center}

Better than either, on every metric, from adding two reciprocals. The gain over
BM25 is about $+0.034$ NDCG@10 --- larger than anything else in this lecture
will buy you. And it is not mysterious: 68 relevant abstracts were found by
exactly one system, and fusion is how you keep both.

\subsubsection{Why not just add the scores?}

The obvious fusion is $\alpha\,\text{BM25} + (1-\alpha)\,\text{cosine}$. It
works, and it costs two things. The scales are incomparable and
corpus-dependent, so the scores must be normalised and the normalisation is
another choice; and $\alpha$ must be tuned on held-out queries, so you now need
labels for a method whose selling point was needing none. Rank fusion avoids
both. Score fusion can do better when you have the labels to tune it --- which
is exactly the condition under which you would have fine-tuned the encoder
instead.

\subsection{Cross-encoders}

\[
  \text{score}(q,d) = \text{Transformer}\big(\,[\text{CLS}]\;q\;[\text{SEP}]\;d\,\big)
\]

Concatenate the query and the document and put both through the model
\emph{together}. Every token of the query can attend to every token of the
document, which is exactly what a bi-encoder forbids by construction.

\begin{center}
\small
\begin{tabular}{p{0.2\textwidth}p{0.32\textwidth}p{0.33\textwidth}}
\toprule
 & \textbf{Bi-encoder} & \textbf{Cross-encoder} \\
\midrule
Sees & each text alone & the pair, jointly \\
Document vectors & precomputed & impossible --- they depend on the query \\
Cost per query & one encoding $+\ N$ dot products & $N$ transformer passes \\
Accuracy & lower & higher \\
\bottomrule
\end{tabular}
\end{center}

\begin{keybox}[title={One sentence worth being able to say}]
The accuracy difference and the cost difference have the \emph{same cause}:
whether the query is present when the document is read.
\end{keybox}

\subsubsection{Why the accurate model cannot go first}

A full cross-encoder scan of this corpus, for these claims, is
\[ 5{,}183 \times 300 = 1{,}554{,}900 \ \text{transformer passes} \]
on the toy corpus of this course. A web index is $10^{10}$ documents and the
user is waiting 200~ms, so the arithmetic is not close --- it is off by ten
orders of magnitude. Hence the pipeline: a cheap first stage reduces millions
to a hundred, and an expensive second stage reorders that hundred. The
architecture of every production search system is a direct consequence of the
fact that the good model does not scale.

\subsection{Re-ranking, measured}

Cross-encoder over BM25's top $k$, everything below $k$ left where it was:

\begin{center}
\small
\begin{tabular}{lrrr}
\toprule
\textbf{System} & \textbf{NDCG@10} & \textbf{P@1} & \textbf{Pairs scored} \\
\midrule
BM25 alone         & 0.6611 & 0.5433 & 0 \\
$+$ re-rank top 10  & 0.6656 & 0.5600 & 3{,}000 \\
$+$ re-rank top 50  & 0.6773 & 0.5667 & 15{,}000 \\
$+$ re-rank top 100 & 0.6781 & 0.5667 & 30{,}000 \\
\bottomrule
\end{tabular}
\end{center}

Going from 10 to 50 buys $+0.0117$; going from 50 to 100 buys $+0.0008$ for
twice the compute. The first ten candidates give most of the gain and the next
ninety give a tenth of it. That is the shape of every re-ranking curve you will
meet, and the reason production depths are usually in the tens rather than the
thousands.

Note also that recall@10 barely moves: re-ranking the top 10 cannot change
\emph{which} documents are in the top 10. It changes their order, and P@1 with
it.

\subsection{The ceiling the re-ranker cannot pass}

Re-ranking reorders; it cannot retrieve. So the recall of the first stage at
depth $k$ is a hard bound on the second stage at depth $k$:

\begin{center}
\small
\begin{tabular}{lrrrr}
\toprule
\textbf{Depth $k$}     & 10 & 50 & 100 & 1000 \\
BM25 recall@$k$        & 0.7784 & 0.8654 & 0.8852 & 0.9650 \\
\bottomrule
\end{tabular}
\end{center}

Re-rank BM25's top 100 and a \emph{perfect} cross-encoder still cannot exceed
recall 0.8852 at any cutoff --- 11.5\% of relevant abstracts are simply not in
the candidate set. Notice that the hybrid first stage reaches recall@100 of
0.9617.

\begin{keybox}[title={The most useful sentence in this lecture}]
Improving the first stage raises the ceiling for everything above it.
Improving the second stage never does.
\end{keybox}

\subsubsection{An exam-style question, worked}

\emph{``A system reports NDCG@10 of 0.6781 using a cross-encoder over the top
100 of a BM25 first stage. Its authors propose replacing the cross-encoder with
a larger one. What is the most you can say about the gain, and what would you
propose instead?''}

\begin{itemize}[leftmargin=*,nosep]
  \item The first stage's recall@100 is 0.8852, so no re-ranker can exceed that
        recall: 11.5\% of relevant documents are absent from the candidates.
  \item Going from depth 50 to 100 already bought only $+0.0008$, so the
        candidates are close to exhausted.
  \item Propose improving the \emph{first} stage. A hybrid one reaches
        recall@100 of 0.9617, which raises the ceiling itself.
\end{itemize}

Section~B of the paper is exactly this shape: a table, a proposal, and the
question of whether the proposal can possibly work.

\section{Approximate nearest neighbours}\label{sec:ann}

The dense scan touches every one of the 8.0~MB of vectors for every query, and
unlike the inverted index it cannot skip anything. The standard answer, and the
one built by hand in the notebook: cluster the vectors once ($k$-means, 64
clusters), then at query time score only the nearest few clusters.

\begin{keybox}[title={What this changes about the numbers}]
Every dense number so far was an exact scan, so it measured the \emph{model}.
From here on the numbers measure the model \emph{and the index}, and the two
must be reported separately --- which is exactly why this whole lecture was
built on a corpus small enough to scan exactly.
\end{keybox}

The claim behind such an index is that a query's nearest documents are mostly
in the clusters nearest the query. That is a claim about geometry and it is
only approximately true: a document near a cluster boundary is missed when its
cluster is not probed, and boundary documents are common in high dimensions. So
the index has a \emph{recall}, and that recall is a parameter rather than a
property.

\begin{failbox}[title={Measure the index against the exact scan}]
Recall of the approximate index is measured against the \emph{exact scan},
never against the relevance judgements. Mixing the two conflates ``my index
missed it'' with ``my model ranked it low'', and those have different fixes.
\end{failbox}

\begin{center}
\small
\begin{tabular}{rrr}
\toprule
\textbf{Clusters probed} & \textbf{Corpus scanned} &
\textbf{Recall@10 vs exact} \\
\midrule
1  & 1.8\%   & 0.5047 \\
2  & 3.6\%   & 0.6580 \\
4  & 7.3\%   & 0.8013 \\
8  & 14.1\%  & 0.8920 \\
16 & 27.9\%  & 0.9577 \\
64 & 100.0\% & 1.0000 \\
\bottomrule
\end{tabular}
\end{center}

Probing 4 of 64 scans 7.3\% of the corpus and keeps 80.1\% of the exact top
ten. A dial, not a default --- and a vector database ships with it set by
somebody who never saw your corpus.

\subsection{A latency budget}

Suppose 200~ms end to end, with the pipeline as built:

\begin{center}
\small
\begin{tabular}{lrl}
\toprule
\textbf{Stage} & \textbf{Model passes} & \textbf{Scales with} \\
\midrule
BM25                            & 0  & query terms $\times$ postings length \\
Encode the query                & 1  & nothing --- the query is short \\
Dense scan (exact)              & 0  & corpus size $\times$ 384 \\
Dense scan (probe 4 of 64)      & 0  & 7.3\% of that \\
Re-rank top 50                  & 50 & the depth you chose \\
\bottomrule
\end{tabular}
\end{center}

Two of those five rows are dials you set, and both trade quality for time along
a curve you can measure. The budget does not tell you where to sit on them; it
does tell you that the choice is yours and has to be made deliberately.

\section{The whole table, in one place}

\begin{center}
\small
\begin{tabular}{lrrl}
\toprule
\textbf{System} & \textbf{NDCG@10} & \textbf{R@100} & \textbf{Cost per query} \\
\midrule
Corpus unranked            & 0.0012 & 0.0117 & nothing \\
Raw term counting          & 0.0944 & ---    & postings walk \\
BM25                       & 0.6611 & 0.8852 & postings walk \\
Dense, zero shot           & 0.6451 & 0.9250 & 1 encode $+$ full scan \\
BM25 $+$ cross-encoder@100 & 0.6781 & 0.8852 & $+$ 100 transformer passes \\
Hybrid (BM25 $+$ dense, RRF) & 0.6948 & 0.9617 & 1 encode $+$ both \\
\bottomrule
\end{tabular}
\end{center}

The best NDCG@10 in this lecture comes from the row with \emph{no training in
it at all} --- two untrained retrievers and a sum of reciprocals.

\subsection{The arc of the two IR lectures}

\begin{center}
\small
\begin{tabular}{p{0.28\textwidth}p{0.62\textwidth}}
\toprule
\textbf{Step} & \textbf{Because} \\
\midrule
Boolean retrieval    & obvious --- and returns nothing for 98\% of claims \\
Score, do not filter & a missing term should cost something, not everything \\
idf, saturation, length & three named failures of raw counting \\
Dense retrieval      & strings cannot capture paraphrase \\
Hybrid               & the two fail on different documents \\
Re-ranking           & the accurate model does not scale \\
Approximate index    & the dense scan cannot skip anything \\
\bottomrule
\end{tabular}
\end{center}

Every row is a response to a \emph{measured} failure of the row above.

\subsection{How to choose, in practice}

\begin{center}
\small
\begin{tabular}{p{0.42\textwidth}p{0.48\textwidth}}
\toprule
\textbf{Situation} & \textbf{What to do} \\
\midrule
No labels, technical vocabulary & BM25. Measure before considering anything
else \\
No labels, everyday language, paraphrase matters & hybrid --- BM25 plus an
off-the-shelf encoder, fused \\
Thousands of labelled pairs in your domain & fine-tune a bi-encoder; now it can
beat BM25 \\
Latency budget allows a second stage & add a cross-encoder over the top
50--100 \\
Corpus too large to scan & approximate index, and report the recall you chose \\
\bottomrule
\end{tabular}
\end{center}

Every row starts with BM25 in the table. Not because it wins --- on your corpus
it may not --- but because a comparison without it is not a measurement.

\begin{trybox}[title={Exercise 20.1}]
Re-rank the \emph{hybrid} top 100 rather than BM25's. The ceiling argument
above predicts what should happen to recall; check whether NDCG follows, and
say why the two need not move together.
\end{trybox}

\section{Reference}

\subsection{Five questions for a neural retrieval claim}

\begin{enumerate}[leftmargin=*]
  \item \term{Is BM25 in the table?} On the same corpus, queries and
        judgements, tuned no less carefully than the neural model.
  \item \term{Zero-shot or fine-tuned?} And if fine-tuned, on what, and how
        close is it to the test corpus?
  \item \term{Exact scan or approximate index?} If approximate, at what recall
        against the exact scan?
  \item \term{Where did the negatives come from?} And what fraction of them are
        unjudged rather than judged irrelevant?
  \item \term{First stage or whole pipeline?} A re-ranker's NDCG is a property
        of the candidates it was given.
\end{enumerate}

\subsection{Where students lose marks}

\begin{itemize}[leftmargin=*,nosep]
  \item Saying a cross-encoder is ``a bigger bi-encoder''. It is a different
        computation: joint, and therefore unindexable.
  \item Claiming a re-ranker improves recall@100. It cannot --- it reorders a
        fixed candidate set.
  \item Reporting approximate-index recall against the judgements rather than
        against the exact scan, which conflates two different errors.
  \item Describing in-batch negatives as ``free''. The scores are free; the
        encodings are not, and the batch is the memory limit.
  \item Concluding that dense retrieval does not work. The claim is narrower:
        \emph{zero-shot}, on \emph{this} corpus.
\end{itemize}

\subsection{Vocabulary}

\begin{center}
\small
\begin{tabular}{ll}
\toprule
Bi-encoder            & query and document encoded separately; score is a dot product \\
Cross-encoder         & the pair encoded jointly; cannot be precomputed \\
In-batch negatives    & the other documents in the batch, used as negatives \\
Hard negative         & a high-scoring irrelevant document, mined from a first stage \\
First stage / re-ranker & cheap over everything; expensive over a few \\
RRF                   & fusing rankings by summing $1/(k+\text{rank})$ \\
ANN, IVF, nprobe      & approximate search; cluster the vectors, probe a few \\
Zero shot             & applied without training on this task or corpus \\
\bottomrule
\end{tabular}
\end{center}

\subsection{Scope}

\term{Examinable.} Vocabulary mismatch and why lexical repairs are bounded;
what a bi-encoder computes and why it is indexable; the contrastive objective
and in-batch negatives; why hard negatives matter and how they go wrong;
cross-encoders and the first-stage/re-rank argument including the ceiling;
hybrid fusion; the recall/latency trade-off of an approximate index.

\term{Not examinable --- engineering.} Specific model checkpoints,
vector-database APIs, HNSW and product-quantisation internals.

\term{Beyond the syllabus.} Late interaction (ColBERT), learned sparse
retrieval (SPLADE), distillation of cross-encoders into bi-encoders. Each is a
direct response to a trade-off this lecture measured, and they are where to
read next.

\subsection{Reading}

\begin{itemize}[leftmargin=*,nosep]
  \item These notes are the primary source.
  \item Karpukhin et al., \emph{Dense Passage Retrieval} --- where in-batch
        negatives were made to work.
  \item Thakur et al., \emph{BEIR} --- the benchmark built to test zero-shot
        generalisation, and the source of the result we reproduced.
  \item Nogueira and Cho, \emph{Passage Re-ranking with BERT} --- the
        cross-encoder as a second stage.
  \item Johnson et al., \emph{Billion-scale similarity search} --- what an IVF
        index is, done properly.
\end{itemize}

\subsection{Summary}

\begin{keybox}[title={One line to keep}]
The first stage sets the ceiling; the second stage only decides how close you
get to it.
\end{keybox}

A lexical index matches strings and relevance is about meaning, and no patch on
the index closes that gap. A bi-encoder can be indexed because $E(d)$ does not
depend on the query --- and that same independence is why it is the weaker
model. Zero-shot on this corpus it loses to BM25 on NDCG@10, 0.6451 against
0.6611, and wins on recall@100, 0.9250 against 0.8852. The two fail on
different documents: 68 of 339 relevant abstracts were found by exactly one, so
fusing them beats both at 0.6948. And the accurate model goes second because it
does not scale --- and it can never exceed the first stage's recall at that
depth.

\medskip

If a lecture in this part has a single practical recommendation, it is this:
\emph{before you replace a lexical retriever, try keeping it.}

\end{document}
